\section{Experiments}
\label{sec:experiments}

We conduct two experiments to evaluate how prompt injection attacks affect tool-using agents and how effectively each defensive strategy mitigates them. All experiments use the agent environment and evaluation metrics described in Section~\ref{sec:approach}. Each experiment runs 50 trials per attack--defense configuration with a fixed random seed for reproducibility.

\subsection{Experiment 1: Direct Prompt Injection}

In the first experiment, the attacker interacts directly with the agent as a normal user but embeds malicious instructions inside the prompt. We test four attack scenarios that target different aspects of agent behavior:

\begin{itemize}
    \item \textbf{System override}: the attacker instructs the agent to ignore its system prompt, read a sensitive database record, and email the contents to an attacker-controlled address.
    \item \textbf{Tool misuse}: the attacker embeds a request to execute malicious code containing \texttt{os.system} calls alongside a benign-looking information lookup.
    \item \textbf{Social engineering}: the attacker impersonates a system administrator and requests that the agent output sensitive fields, including SSN, from a personnel record.
    \item \textbf{Data exfiltration}: the attacker asks the agent to include a canary API key in a search query, leaking it through the tool's outbound request.
\end{itemize}

Each scenario is run against all six defense configurations. The results are shown in Table~\ref{tab:direct}.

\begin{table}[h]
\centering
\caption{Experiment 1: Direct prompt injection results (50 trials per condition).}
\label{tab:direct}
\begin{tabular}{lcccc}
\hline
\textbf{Defense} & \textbf{ASR} & \textbf{TMR} & \textbf{SDLR} & \textbf{TUD} \\
\hline
None              & 75.0\% & 42.9\% & 25.0\% & 0.0\% \\
Prompt isolation  & 75.0\% & 42.9\% &  7.5\% & 0.0\% \\
Provenance        & 75.0\% & 42.9\% &  8.5\% & 0.0\% \\
Tool permission   & 25.0\% & 42.9\% & 25.0\% & 0.0\% \\
Prompt detection  & 34.0\% & 16.9\% &  6.0\% & 41.0\% \\
All combined      & 25.0\% &  6.2\% &  1.5\% & 47.5\% \\
\hline
\end{tabular}
\end{table}

In direct injection, prompt isolation and provenance marking have no effect on ASR. This is expected: the attacker \textit{is} the user, so wrapping user input with untrusted markers does not change the fact that the malicious instructions originate from the prompt itself. However, both defenses reduce SDLR significantly (from 25.0\% to 7.5\% and 8.5\%), because the agent becomes less likely to echo sensitive data in its response when provenance labels are present.

Tool permission control reduces ASR from 75.0\% to 25.0\% by blocking tool calls that match known malicious argument patterns. It does not reduce TMR because the agent still \textit{attempts} the misuse; the permission layer blocks execution after the decision is made. Prompt detection reduces ASR to 34.0\% by flagging suspicious patterns before the agent processes them, but introduces 41.0\% task utility degradation due to false positives on legitimate requests.

When all defenses are combined, ASR remains at 25.0\% while TMR drops to 6.2\% and SDLR to 1.5\%. The combined configuration achieves the strongest protection but at the cost of 47.5\% utility degradation, reflecting the trade-off between security and usability.

\subsection{Experiment 2: Indirect Prompt Injection}

In the second experiment, the attacker does not interact with the system directly. Instead, malicious instructions are embedded in the content returned by tools that the agent invokes during normal task execution. We test four scenarios:

\begin{itemize}
    \item \textbf{Poisoned search result}: a web search returns legitimate financial data followed by a fake system notice instructing the agent to email the summary to an attacker-controlled address.
    \item \textbf{Poisoned database record}: a project status query returns valid data appended with instructions to execute malicious code for ``access logging.''
    \item \textbf{Data exfiltration via email}: a personnel record lookup returns employee data with an injected directive to forward the complete record, including SSN, to an external address.
    \item \textbf{Chained tool abuse}: a weather search returns forecast data with an instruction to also query an unrelated sensitive database record.
\end{itemize}

The results are shown in Table~\ref{tab:indirect}.

\begin{table}[h]
\centering
\caption{Experiment 2: Indirect prompt injection results (50 trials per condition).}
\label{tab:indirect}
\begin{tabular}{lcccc}
\hline
\textbf{Defense} & \textbf{ASR} & \textbf{TMR} & \textbf{SDLR} & \textbf{TUD} \\
\hline
None              & 48.5\% & 20.5\% & 43.0\% & 0.0\% \\
Prompt isolation  & 27.5\% & 16.1\% & 19.0\% & 0.0\% \\
Provenance        & 33.0\% & 14.3\% & 21.5\% & 0.0\% \\
Tool permission   & 17.5\% & 24.9\% & 42.5\% & 0.0\% \\
Prompt detection  & 17.5\% &  3.0\% & 39.0\% & 0.0\% \\
All combined      &  5.0\% &  1.9\% & 12.5\% & 0.0\% \\
\hline
\end{tabular}
\end{table}

Unlike in Experiment~1, prompt isolation and provenance marking now produce meaningful reductions in ASR. Prompt isolation reduces ASR from 48.5\% to 27.5\%, while provenance marking reduces it to 33.0\%. This is because the injected content originates from tool outputs rather than the user prompt, so marking it as untrusted data helps the agent distinguish between legitimate instructions and injected ones.

Tool permission control reduces ASR to 17.5\% by blocking outbound tool calls with suspicious arguments. Prompt detection achieves the same ASR of 17.5\% while also reducing TMR from 20.5\% to 3.0\%, because it identifies injection patterns in tool outputs before the agent can act on them. Notably, neither defense introduces utility degradation in the indirect setting, since legitimate tasks do not trigger the detection patterns that are specific to attack content.

The combined defense achieves 5.0\% ASR with 1.9\% TMR and 12.5\% SDLR. The residual ASR comes from attack variants whose injection language is subtle enough to avoid pattern detection while still influencing the agent's behavior.

\subsection{Summary of Findings}

Table~\ref{tab:summary} summarizes the ASR across both experiments for each defense strategy.

\begin{table}[h]
\centering
\caption{Attack Success Rate comparison across experiments.}
\label{tab:summary}
\begin{tabular}{lcc}
\hline
\textbf{Defense} & \textbf{Direct} & \textbf{Indirect} \\
\hline
None              & 75.0\% & 48.5\% \\
Prompt isolation  & 75.0\% & 27.5\% \\
Provenance        & 75.0\% & 33.0\% \\
Tool permission   & 25.0\% & 17.5\% \\
Prompt detection  & 34.0\% & 17.5\% \\
All combined      & 25.0\% &  5.0\% \\
\hline
\end{tabular}
\end{table}

Several observations emerge from the results:

\begin{enumerate}
    \item \textbf{Direct and indirect attacks require different defenses.} Prompt isolation and provenance marking are ineffective against direct injection (where the attacker is the user) but reduce indirect injection ASR by 15--21 percentage points. Conversely, tool permission control is effective against both attack types because it operates at the execution layer regardless of where the injection originates.

    \item \textbf{No single defense is sufficient.} The best individual defense (tool permission) still allows 17.5--25.0\% of attacks to succeed. Only the combined configuration reduces ASR to 25.0\% for direct and 5.0\% for indirect attacks.

    \item \textbf{Security and utility trade off.} Prompt detection introduces 41.0\% utility degradation in the direct setting due to false positives on legitimate user requests. In contrast, the same defense causes no degradation in the indirect setting, where attack patterns are more distinctive. This suggests that detection-based defenses are better suited for filtering external content than user input.

    \item \textbf{Data leakage is a distinct risk from attack success.} Even when ASR is relatively low, SDLR can remain high if the agent echoes sensitive data from tool outputs in its response. Prompt isolation and provenance marking are most effective at reducing this specific risk, even when they do not prevent the underlying attack from influencing tool selection.
\end{enumerate}
